\documentclass{jsarticle}
\usepackage{amsmath}
\usepackage{amssymb}
\begin{document}
\title{空間概念の量子化は、光量子仮設を使う}
\author{Masaaki Yamaguchi}
\date{}
\maketitle


    空間概念の量子化のため、最低エネルギーのエントロピー値の式が、
    $$ e^{f} + e^{-f} \ge e^{f} - e^{-f} $$
    として、この$e^{-x \log x}$が、光量子仮設を説明している。
    量子飛躍は、
    $$ h\nu = [e^{-x \log x}]$$となる。
    このため、飛び飛びの値をとる。
    最低エネルギーのエントロピー値$e^{-x \log x}$のため、
    $$ [\eta + \bar{h}(x)]dx^{\mu}dx^{\nu} $$
    $$ e^{-2\pi T||\psi||} [\eta + \bar{h}(x)]dx^{\mu}dx^{\nu} $$
    これは、単体量であり、
    $$T^2d^2 \psi$$
    これは、アーベル多様体による熱エネルギーを表している。
    ブラックホールのエネルギー体である。
    この単体量と熱エネルギーが同等だと、
    $$ = ||ds^2||$$
    平坦な宇宙の場合となる。
    空洞放射は、無限の熱エネルギーなのは、$e^{x \log x}$
    ヒッグス場のエネルギーより
    $$ {d \over df}F = m(x) = e^{x \log x} $$
    ここで、ガンマ関数とオイラーの定数の多様体積分が、
    $$ \int \Gamma(\gamma)^{'}dx_m = e^{-x \log x}$$
    $$ \int Cdx_m = e^{-x \log x} $$
    のために、光電効果から、フェルマーの定理の
    $$ e^{-x \log x} \le z^n \le e^{x \log x}$$
    としての範囲を、単体量でのエネルギー範囲をもつ。
    xの値をとるために、飛び飛びの値を取るのと、
    無限のエネルギーを持つのを説明できる。
    全ては、ゼータ関数から求まる。
    　


\end{document}
